\chapter{Tinjauan Pustaka dan Dasar Teori}

\section{Tinjauan Pustaka}

Kajian penerapan \textit{blockchain} untuk mata uang digital bank sentral
mencakup beragam fokus. Bagian ini meninjau penelitian terdahulu yang relevan
dengan rancang bangun sistem Rupiah Digital ritel berbasis Hyperledger Fabric,
meliputi arsitektur, transfer ritel, pengawasan, dan evaluasi kinerja.

\subsection{Penelitian Terdahulu tentang Desain dan Privasi CBDC}

Sethaput dan Innet \cite{sethaput2025} menyusun \textit{systematization of
knowledge} (SoK) mengenai aplikasi \textit{blockchain} pada CBDC. Studi tersebut
mengidentifikasi pilihan antara \textit{token-based} dan \textit{account-based},
model satu jenjang dan dua jenjang, serta tingkat privasi sebagai dimensi desain
yang paling menentukan. Lukito dkk. \cite{lukito2026governance} membandingkan
sembilan program CBDC berdasarkan dokumen publik dan menguraikan perbedaan peran
perantara dalam distribusi dua jenjang.

Bhawana dan Kumar \cite{bhawana2021} membangun jaringan CBDC berbasis
\textit{permissioned blockchain}. Menurut studi tersebut, keanggotaan terbatas
memberikan kendali tata kelola yang dibutuhkan bank sentral tanpa mengorbankan
transparansi. Islam dan In \cite{islam2023} menggunakan \textit{consortium
blockchain} dengan model \textit{unspent transaction output} (UTXO) untuk menjaga
privasi sekaligus auditabilitas. Pocher dan Veneris \cite{Pocher2022Privacy}
membahas privasi dan transparansi CBDC melalui skema \textit{anti-money
laundering} dan \textit{countering the financing of terrorism} (AML/CFT) berbasis
regulasi-oleh-desain. Ketiga studi tersebut tidak berfokus pada Rupiah Digital.

Aryani \cite{aryani2023critlit} melakukan tinjauan literatur kritis terhadap desain
CBDC Indonesia dan mengusulkan arsitektur dua jenjang Digital Rupiah. Arsitektur
tersebut memuat lapisan \textit{Digital Rupiah Depository} sebagai perantara antara
Bank Indonesia dan peserta jaringan, serta koneksi ke BI-RTGS dan jaringan
\textit{blockchain} lintas batas. Penelitian ini mengadopsi prinsip distribusi
berjenjang tersebut dengan dua penyederhanaan yang disengaja sesuai cakupan
purwarupa. Pertama, lapisan \textit{Digital Rupiah Depository} diciutkan menjadi
satu organisasi distributor (\texttt{DistributorMSP}) yang menerima penerbitan
langsung dari \texttt{CentralBankMSP}; pemisahan \textit{wholesale} dan
\textit{retail} dilakukan sebagai jenis bank di dalam organisasi tersebut. Kedua,
integrasi BI-RTGS dan jaringan lintas batas tidak diimplementasikan dan dinyatakan
sebagai batasan cakupan penelitian (Bab~I).

\subsection{Penelitian Terdahulu Berbasis Hyperledger Fabric}

Rohmalia dan Nugraha \cite{rohmalia2023retail} menyusun model CBDC ritel Indonesia
berbasis Hyperledger Fabric yang mencakup pembuatan akun, konversi aset, dan
transfer antarpengguna. Hasil implementasi yang mereka laporkan memperlihatkan
pembuatan akun dan transfer, disertai waktu pengiriman untuk satu dan 50 data.
Makalah tersebut tidak menyajikan pengawasan multipihak ataupun tolok ukur Caliper
dengan artefak yang dapat diperiksa.

Kuzlu dkk. \cite{fabric_perf2023} menganalisis kinerja kerangka Hyperledger Fabric
pada aspek \textit{throughput}, latensi, dan skalabilitas. Studi tersebut berfokus
pada kerangka Fabric secara umum, sedangkan skenario CBDC ritel tidak menjadi
fokusnya. Angka kinerjanya tidak dibandingkan langsung dengan penelitian ini karena
versi, topologi, dan beban kerja berbeda.

\subsection{Posisi dan Celah Penelitian}

Tabel~\ref{tab:ch2-tinjauan} merangkum posisi kedelapan penelitian tersebut
terhadap penelitian ini.

\begin{table}[htbp]
  \caption{Perbandingan penelitian terdahulu dengan penelitian ini.}
  \vspace{0.5em}
  \centering
  \begingroup
  \renewcommand{\arraystretch}{1.25}
  \setlength{\tabcolsep}{3pt}
  \small
  \begin{tabular}{>{\RaggedRight\arraybackslash}p{0.16\textwidth}
                  >{\RaggedRight\arraybackslash}p{0.23\textwidth}
                  >{\RaggedRight\arraybackslash}p{0.29\textwidth}
                  >{\RaggedRight\arraybackslash}p{0.25\textwidth}}
    \toprule
    \textbf{Penulis/Tahun} & \textbf{Metode} & \textbf{Kontribusi} & \textbf{Keterbatasan} \\
    \midrule
    Sethaput dan Innet (2025)~\cite{sethaput2025} &
    SoK (\textit{systematization of knowledge}) &
    Dimensi desain CBDC berbasis \textit{blockchain}: \textit{token-based} dan \textit{account-based}, satu dan dua jenjang, serta privasi &
    Sistematisasi literatur tanpa pengujian purwarupa \\
    Lukito dkk. (2026)~\cite{lukito2026governance} &
    Perbandingan sembilan program CBDC berdasarkan dokumen publik &
    Perbedaan peran perantara pada distribusi dua jenjang &
    Bergantung pada dokumen publik; tidak menguji implementasi \\
    Bhawana dan Kumar (2021)~\cite{bhawana2021} &
    Rancangan jaringan CBDC pada \textit{permissioned blockchain} &
    Kendali tata kelola bank sentral tanpa mengorbankan transparansi &
    Fitur ritel tidak terintegrasi; tolok ukur kinerja terbatas \\
    Rohmalia dan Nugraha (2023)~\cite{rohmalia2023retail} &
    Model CBDC ritel dan implementasi pada Hyperledger Fabric &
    Pembuatan akun, konversi aset, dan transfer antarpengguna &
    Tanpa pengawasan multipihak dan tolok ukur Caliper yang dapat diperiksa \\
    Aryani (2023)~\cite{aryani2023critlit} &
    Tinjauan literatur kritis &
    Arsitektur dua jenjang Digital Rupiah dengan \textit{Digital Rupiah Depository} &
    Bersifat konseptual tanpa purwarupa \\
    Islam dan In (2023)~\cite{islam2023} &
    Rancangan CBDC pada \textit{consortium blockchain} dengan model UTXO &
    Privasi sekaligus auditabilitas &
    Tidak mencakup fitur ritel; tanpa tolok ukur kinerja \\
    Pocher dan Veneris (2022)~\cite{Pocher2022Privacy} &
    Skema AML/CFT berbasis regulasi-oleh-desain &
    Privasi dan transparansi CBDC melalui skema AML/CFT &
    Tidak berfokus pada Rupiah Digital maupun Proyek Garuda \\
    Kuzlu dkk. (2019)~\cite{fabric_perf2023} &
    Analisis kinerja kerangka \textit{Hyperledger Fabric} &
    \textit{Throughput}, latensi, dan skalabilitas sebagai rujukan kinerja Fabric &
    Berfokus pada kerangka Fabric secara umum, bukan skenario CBDC ritel \\
    \midrule
    \textbf{Penelitian ini} &
    Purwarupa Fabric dengan pengujian fungsional dan Caliper &
    Penerbitan dua jenjang, transfer ritel, pengawasan, dan protokol Caliper dalam satu artefak &
    Privasi transaksi, arsitektur hibrida, dan interkoneksi lintas batas tidak diimplementasikan; satu \textit{application channel} dan satu \textit{orderer} \\
    \bottomrule
  \end{tabular}
  \endgroup
  \label{tab:ch2-tinjauan}
\end{table}

\textbf{Kebaruan.} Dalam batas korpus yang ditinjau, kontribusi penelitian ini
terletak pada integrasi alur penerbitan dua jenjang, transaksi ritel, pengawasan,
serta protokol Caliper dalam satu artefak Fabric yang dapat diperiksa. Kontribusi
ini bukan klaim kebaruan arsitektur CBDC secara universal maupun klaim bahwa
platform yang digunakan Proyek Garuda adalah Fabric \cite{ree2023enaira,
bi_garuda_poc}. Keterulangan didukung oleh manifes yang mencatat commit dan hash
masukan untuk setiap batch benchmark.

\textbf{Celah penelitian.} Dalam korpus yang ditinjau, dokumentasi publik Proyek
Garuda baru menjangkau lapisan \textit{wholesale} \cite{bi_garuda_poc}, sedangkan
Rupiah Digital ritel ditempatkan sebagai tahap lanjutan dalam peta jalan dan
dokumen konsultatif \cite{bi_whitepaper2022, bi_consultative}. Tidak ada rancangan
publik dalam korpus tersebut yang memperlihatkan kontrol \textit{know your customer}
(KYC) dan AML/CFT ritel yang ditegakkan dan dapat diverifikasi pada buku besar.
Penelitian ini menjawab celah tersebut melalui dua pertanyaan penelitian. Pertama,
bagaimana purwarupa Rupiah Digital ritel di atas Hyperledger Fabric dirancang dan
dibangun dengan kontrol KYC dan AML/CFT terpilih, serta bagaimana kontrol tersebut
diverifikasi secara fungsional? Kedua, bagaimana kinerja purwarupa tersebut diukur
dengan Hyperledger Caliper?

\section{Dasar Teori}

\subsection{Central Bank Digital Currency}

\textit{Central Bank Digital Currency} (CBDC) adalah representasi digital dari
mata uang suatu negara yang diterbitkan oleh dan menjadi kewajiban langsung bank
sentral \cite{unknown2024p}. Sebagai kewajiban bank sentral, CBDC memiliki
kualitas bebas risiko yang setara dengan uang tunai, berbeda dengan uang
elektronik yang merupakan kewajiban penerbit swasta. Secara umum CBDC dibedakan
menjadi dua jenis berdasarkan pengguna sasarannya, yaitu CBDC \textit{wholesale}
yang ditujukan untuk transaksi antarlembaga keuangan, dan CBDC \textit{retail}
yang ditujukan untuk masyarakat luas \cite{unknown2022f,
SyarifuddinBakhtiar2021Monetary}.

Dalam penerapannya, sebagian besar bank sentral memilih model dua jenjang
(\textit{two-tier}), yaitu bank sentral menerbitkan CBDC dan mengelola
\textit{ledger} pada tingkat \textit{wholesale}, sementara distribusi kepada
masyarakat dilakukan melalui perantara seperti bank. Model ini menjaga peran intermediasi lembaga keuangan dan mengurangi
risiko \textit{disintermediation} pada sistem perbankan \cite{sethaput2025}. Di
Indonesia, Bank Indonesia mengembangkan CBDC dengan nama Rupiah Digital melalui
Proyek Garuda, yang menempuh pendekatan bertahap mulai dari Rupiah Digital
\textit{wholesale} kemudian diperluas menuju Rupiah Digital \textit{retail}
\cite{bi_whitepaper2022}. Desain Rupiah Digital menekankan interoperabilitas
dengan infrastruktur sistem pembayaran dan setelmen, serta penerapan model dua
jenjang dengan pihak perantara
\cite{bi_consultative, bi_garuda_poc}. Taksonomi CBDC berdasarkan segmen
pengguna dan model distribusi ditunjukkan pada Gambar~\ref{fig:cbdc-taxonomy}.

% Taksonomi CBDC berdasarkan segmen pengguna dan model distribusi (TikZ).
% Di-input dari BAB II §2.2.1.

\begin{figure}[H]
\centering
\resizebox{0.88\textwidth}{!}{%
\begin{tikzpicture}[
  root/.style={draw=black, fill=gray!30, rectangle, rounded corners=3pt,
    minimum width=2.8cm, minimum height=0.9cm, align=center, font=\small\bfseries},
  node1/.style={draw=black, fill=gray!12, rectangle, rounded corners=3pt,
    minimum width=3.0cm, minimum height=0.9cm, align=center, font=\footnotesize\bfseries},
  node2/.style={draw=black, fill=white, rectangle, rounded corners=3pt,
    minimum width=3.8cm, minimum height=1.6cm, align=center, font=\footnotesize},
  selected/.style={draw=black, fill=black!15, rectangle, rounded corners=3pt,
    minimum width=3.8cm, minimum height=1.6cm, align=center, font=\footnotesize\bfseries},
  ar/.style={-{Stealth}, thick},
]

%--- Akar ---
\node[root] (cbdc) at (5.5, 4.5) {CBDC};

%--- Level 1 ---
\node[node1] (wh) at (2.2, 3.0)  {Wholesale\\(Antarlembaga)};
\node[node1] (rt) at (8.8, 3.0)  {Retail\\(Masyarakat Umum)};

%--- Level 2 (Retail sub-types) ---
\node[node2] (one) at (6.4, 1.5)
  {Satu Jenjang\\(\textit{Direct})\\Bank Sentral\\$\to$ Publik};
\node[selected] (two) at (11.2, 1.5)
  {Dua Jenjang\\(\textit{Two-Tier})\\BI $\to$ Bank\\$\to$ Publik};

%--- Arrows ---
\draw[ar] (cbdc.south) to[out=-150,in=90] (wh.north);
\draw[ar] (cbdc.south) to[out=-30,in=90]  (rt.north);
\draw[ar] (rt.south)   to[out=-150,in=90] (one.north);
\draw[ar] (rt.south)   to[out=-30,in=90]  (two.north);

%--- Contoh ---
\node[font=\scriptsize, align=center] at (2.2, 1.8)
  {Contoh: PoC Proyek Garuda\\\cite{bi_garuda_poc}};
\node[font=\scriptsize, align=center, anchor=north] at (6.4, 0.45)
  {Contoh: DCash (Karibia)};
\node[font=\scriptsize, align=center, anchor=north] at (11.2, 0.45)
  {Contoh: eNaira, e-CNY,\\Rupiah Digital};

\end{tikzpicture}%
}
\caption{Taksonomi CBDC berdasarkan segmen pengguna dan model distribusi; model dua jenjang yang diadopsi pada penelitian ini ditampilkan dengan latar abu-abu sesuai rancangan Rupiah Digital Bank Indonesia}
\label{fig:cbdc-taxonomy}
\end{figure}


Peta jalan resmi Proyek Garuda menempatkan pengembangan Rupiah Digital dalam
tiga fase, yaitu \textit{immediate state} untuk \textit{w-Rupiah Digital} berbasis
\textit{cash ledger}, \textit{intermediate state} untuk \textit{w-Rupiah Digital}
dengan kas dan sekuritas digital, serta \textit{end state} yang mengintegrasikan
\textit{w-Rupiah Digital} dengan \textit{r-Rupiah Digital}. Pada fase \textit{end state}, use
case \textit{r-Rupiah Digital} mencakup distribusi, penagihan, dan pembayaran P2P, baik
melalui perantara maupun langsung dari bank sentral sebagaimana ditunjukkan
pada Gambar~\ref{fig:bi-whitepaper-roadmap} \cite{bi_whitepaper2022}.

\begin{figure}[H]
    \centering
    \rotatebox{90}{\includegraphics[width=0.78\textheight]{contents/figures/reference/bi-whitepaper-roadmap.png}}
    \caption{Peta jalan pengembangan Rupiah Digital dalam Proyek Garuda. Sumber: Bank Indonesia \cite{bi_whitepaper2022}.}
    \label{fig:bi-whitepaper-roadmap}
\end{figure}

\subsubsection{Model Pembayaran Berbasis Akun dan Token}

Selain dibedakan berdasarkan segmen pengguna dan model distribusi, CBDC juga
dapat diklasifikasikan berdasarkan mekanisme pembayaran berbasis akun
(\textit{account-based}) atau berbasis token (\textit{token-based}). Pada model
berbasis akun, akses pembayaran bergantung pada hubungan antara identitas
pengguna, akun atau dompet, dan catatan saldo yang dikelola sistem. Pada model
berbasis token, validitas pembayaran terutama bergantung pada verifikasi objek
token atau instrumen nilai yang dipindahkan, sehingga penekanan desain berada
pada keaslian token dan hak pemegangnya untuk membelanjakan token tersebut
\cite{BechGarratt2017Central,Auer2020The}. Perbandingan alur verifikasi kedua
model ditunjukkan pada Gambar~\ref{fig:akun-token}.

% Perbandingan model pembayaran berbasis akun dan token (TikZ).
% Diadaptasi dari Auer & Bohme (2020). Di-input dari BAB II.

\begin{figure}[H]
\centering
\resizebox{0.92\textwidth}{!}{%
\begin{tikzpicture}[
  box/.style={draw=black, fill=white, rectangle, rounded corners=2pt,
    minimum width=2.7cm, minimum height=0.9cm, align=center, font=\footnotesize},
  ver/.style={draw=black, fill=gray!15, rectangle, rounded corners=2pt,
    minimum width=2.4cm, minimum height=0.9cm, align=center, font=\footnotesize\bfseries},
  term/.style={draw=black, fill=gray!25, rectangle, rounded corners=8pt,
    minimum width=2.6cm, minimum height=0.8cm, align=center, font=\footnotesize},
  ar/.style={-{Stealth}, thick},
]

%==== (a) Berbasis Akun ====
\node[box]  (a-user) at (0, 0.9)  {Pengguna A};
\node[box]  (a-id)   at (0, -0.7) {Dokumen\\identitas};
\node[ver]  (a-ver)  at (3.6, 0.1) {Verifikasi\\identitas};
\node[term] (a-ok)   at (7.6, 0.9) {Transfer ke\\Pengguna B};
\node[term] (a-no)   at (7.6, -0.9) {Ditolak};
\draw[ar] (a-user.east) -- (a-ver.north west);
\draw[ar] (a-id.east)   -- (a-ver.south west);
\draw[ar] (a-ver.east) -- node[above,font=\scriptsize]{Ya} (a-ok.west);
\draw[ar] (a-ver.east) -- node[below,font=\scriptsize]{Tidak} (a-no.west);
\node[font=\footnotesize\bfseries] at (3.6, -2.0)
  {(a) Berbasis akun (\textit{account-based})};

%==== (b) Berbasis Token ====
\node[box]  (b-priv) at (0, -4.5)   {Pengguna A\\enkripsi\\kunci privat};
\node[box]  (b-pub)  at (3.6, -4.5) {Verifikasi\\kunci publik\\Pengguna A};
\node[ver]  (b-ver)  at (7.2, -4.5) {Verifikasi};
\node[term] (b-ok)   at (10.8, -3.7) {Eksekusi transfer\\ke Pengguna B};
\node[term] (b-no)   at (10.8, -5.5) {Ditolak};
\draw[ar] (b-priv.east) -- (b-pub.west);
\draw[ar] (b-pub.east)  -- (b-ver.west);
\draw[ar] (b-ver.east) -- node[above,font=\scriptsize]{Ya} (b-ok.west);
\draw[ar] (b-ver.east) -- node[below,font=\scriptsize]{Tidak} (b-no.west);
\node[font=\footnotesize\bfseries] at (5.4, -6.6)
  {(b) Berbasis token (\textit{token-based})};

\end{tikzpicture}%
}
\caption{Perbandingan model pembayaran CBDC berbasis akun dan berbasis token: model akun memverifikasi identitas pemegang, sedangkan model token memverifikasi keaslian objek nilai. Diadaptasi dari Auer dan B\"ohme \cite{Auer2020The}}
\label{fig:akun-token}
\end{figure}


White Paper CBDC Bank Indonesia menempatkan Rupiah Digital dalam konfigurasi
desain yang mencakup pemisahan \textit{w-Rupiah Digital} dan \textit{r-Rupiah Digital}, pilihan
satu jenjang atau dua jenjang, akses berbasis token dan akun, serta
infrastruktur DLT \textit{permissioned} sebagaimana ditunjukkan pada
Gambar~\ref{fig:bi-whitepaper-design-configuration}. Konfigurasi tersebut
menjadi dasar pembatasan ruang lingkup penelitian ini: purwarupa difokuskan
pada model ritel dua jenjang berbasis dompet akun di atas jaringan
\textit{permissioned blockchain}, bukan implementasi penuh seluruh konfigurasi
Rupiah Digital \cite{bi_whitepaper2022}.

\begin{figure}[H]
    \centering
    \rotatebox{90}{\includegraphics[width=0.78\textheight]{contents/figures/reference/bi-whitepaper-design-configuration.png}}
    \caption{Konfigurasi desain Rupiah Digital menurut White Paper CBDC Bank Indonesia. Sumber: Bank Indonesia \cite{bi_whitepaper2022}.}
    \label{fig:bi-whitepaper-design-configuration}
\end{figure}

Pembedaan ini relevan karena memengaruhi desain privasi dan auditabilitas
transaksi. Syarifuddin dan Bakhtiar menempatkan
\textit{account-based} dan \textit{token-based payment system} sebagai salah
satu dimensi desain CBDC yang perlu dipertimbangkan dalam strategi kebijakan
moneter \cite{SyarifuddinBakhtiar2021Monetary}. Dalam konteks Indonesia,
dokumen konsultatif Rupiah Digital menyatakan bahwa akses pada
Rupiah Digital \textit{wholesale} menggunakan verifikasi berbasis token karena
dianggap sesuai untuk transaksi pasar keuangan yang lebih kompleks
\cite{bi_consultative}. Namun, dokumen tersebut menempatkan Rupiah Digital
ritel sebagai bagian dari tahap integrasi \textit{end state}; karena itu,
penelitian ini tidak mengklaim bahwa desain pembayaran Rupiah Digital ritel
Bank Indonesia telah final sebagai berbasis akun ataupun berbasis token.

\subsubsection{Arsitektur Penerbitan CBDC}

Selain model distribusi satu atau dua jenjang, penerbitan CBDC dikategorikan
menurut struktur klaim menjadi tiga arsitektur, yaitu langsung
(\textit{direct}), tidak langsung (\textit{indirect}), dan hibrida
(\textit{hybrid}) \cite{Auer2020The}. Pada arsitektur langsung, bank sentral
menangani layanan pembayaran ritel secara penuh; pada arsitektur tidak langsung,
klaim ritel berada pada perantara yang menjamin CBDC; sedangkan pada arsitektur
hibrida, CBDC tetap menjadi klaim langsung bank sentral tetapi perantara
mengelola interaksi dengan pengguna. Syarifuddin dan Bakhtiar menempatkan pilihan
arsitektur ini sebagai salah satu dimensi desain CBDC yang memengaruhi strategi
kebijakan moneter \cite{SyarifuddinBakhtiar2021Monetary}. Ketiga arsitektur
tersebut dibandingkan pada Gambar~\ref{fig:arsitektur-cbdc-tipe}.

% Perbandingan arsitektur penerbitan CBDC: tidak langsung, langsung, hibrida (TikZ).
% Diadaptasi dari Auer & Bohme (2020). Di-input dari BAB II.

\begin{figure}[H]
\centering
\resizebox{0.9\textwidth}{!}{%
\begin{tikzpicture}[
  cb/.style={draw=black, fill=gray!30, rectangle, rounded corners=3pt,
    minimum width=2.4cm, minimum height=1.0cm, align=center, font=\footnotesize\bfseries},
  mid/.style={draw=black, fill=gray!12, rectangle, rounded corners=3pt,
    minimum width=2.6cm, minimum height=1.0cm, align=center, font=\footnotesize},
  usr/.style={draw=black, fill=white, rectangle, rounded corners=3pt,
    minimum width=2.2cm, minimum height=1.0cm, align=center, font=\footnotesize},
  lbl/.style={font=\footnotesize\bfseries, align=left},
  note/.style={font=\scriptsize, align=center},
  ar/.style={-{Stealth}, thick},
  ard/.style={-{Stealth}, thick, dashed},
]

%==== Baris 1: Tidak Langsung (Indirect / dua jenjang) ====
\node[lbl] at (-3.6, 2.2) {Tidak langsung\\(\textit{indirect})};
\node[cb]  (cb1)  at (0, 2.2) {Bank\\Sentral};
\node[mid] (int1) at (4.0, 2.2) {Perantara\\(Bank)};
\node[usr] (u1)   at (8.0, 2.2) {Pengguna};
\draw[ar] (cb1.east)  -- (int1.west);
\draw[ar] (int1.east) -- (u1.west);
\node[note] at (4.0, 1.1) {klaim ritel ke perantara};

%==== Baris 2: Langsung (Direct / satu jenjang) ====
\node[lbl] at (-3.6, 0.0) {Langsung\\(\textit{direct})};
\node[cb]  (cb2) at (0, 0.0) {Bank\\Sentral};
\node[usr] (u2)  at (8.0, 0.0) {Pengguna};
\draw[ar] (cb2.east) -- (u2.west);
\node[note] at (4.0, -0.55) {klaim langsung ke bank sentral; ritel ditangani bank sentral};

%==== Baris 3: Hibrida (Hybrid) ====
\node[lbl] at (-3.6, -2.4) {Hibrida\\(\textit{hybrid})};
\node[cb]  (cb3)  at (0, -2.4) {Bank\\Sentral};
\node[mid] (int3) at (4.0, -2.4) {Bank};
\node[usr] (u3)   at (8.0, -2.4) {Pengguna};
\draw[ar] (cb3.east)  -- (int3.west);
\draw[ar] (int3.east) -- (u3.west);
\draw[ard] (cb3.south) to[out=-40,in=-140] node[note,below]{klaim langsung ke bank sentral} (u3.south);
\node[note] at (4.0, -1.3) {Bank menangani pembayaran ritel};

\end{tikzpicture}%
}
\caption{Perbandingan arsitektur penerbitan CBDC: tidak langsung (\textit{indirect}/dua jenjang), langsung (\textit{direct}/satu jenjang), dan hibrida (\textit{hybrid}). Diadaptasi dari Auer dan B\"ohme \cite{Auer2020The}}
\label{fig:arsitektur-cbdc-tipe}
\end{figure}


\subsection{Blockchain dan Distributed Ledger Technology}

\subsubsection{Struktur Blockchain}

\textit{Distributed Ledger Technology} (DLT) adalah teknologi pencatatan
terdistribusi yang memungkinkan banyak pihak memelihara salinan basis data yang
konsisten tanpa otoritas terpusat tunggal. \textit{Blockchain} merupakan salah
satu bentuk DLT yang menyusun transaksi ke dalam blok-blok yang saling
terhubung secara kriptografis. Setiap blok memuat \textit{hash} dari blok
sebelumnya sehingga membentuk riwayat tambah-saja yang perubahan retrospektifnya
dapat dideteksi; perubahan pada satu blok akan merusak keterkaitan \textit{hash}
pada seluruh blok berikutnya \cite{unknown2022}. Sifat ini mendukung integritas
dan ketertelusuran data selama mekanisme konsensus dan kendali akses tetap
berfungsi. Struktur rantai blok diilustrasikan pada
Gambar~\ref{fig:block-structure}.
Setiap blok juga memuat \textit{Merkle root}, yaitu nilai hash yang dihasilkan
dari pohon Merkle seluruh transaksi dalam blok tersebut, sehingga integritas
setiap transaksi dapat diverifikasi secara efisien tanpa harus membaca semua
data transaksi \cite{unknown2022}. Struktur pohon Merkle ditunjukkan pada
Gambar~\ref{fig:merkle-tree}.

% Struktur rantai blok dan pohon Merkle (TikZ).
% Di-input dari BAB II §2.2.1.1.

\begin{figure}[H]
\centering
\resizebox{\textwidth}{!}{%
\begin{tikzpicture}[
  blkbody/.style={draw=black, fill=white, rectangle, minimum width=3.6cm,
    minimum height=0.5cm, align=center, font=\scriptsize, inner sep=2pt},
  blkhdr/.style={draw=black, fill=gray!20, rectangle, minimum width=3.6cm,
    minimum height=0.6cm, align=center, font=\scriptsize\bfseries, inner sep=2pt},
  ar/.style={-{Stealth}, thick},
]

%--- Blok 0 (Genesis) ---
\node[blkhdr] (b0h)  at (0, 0)     {Blok 0 (Genesis)};
\node[blkbody, below=0pt of b0h] (b0p) {Hash sebelumnya: \texttt{0x000\ldots}};
\node[blkbody, below=0pt of b0p] (b0m) {Merkle Root: \texttt{A1B2}};
\node[blkbody, below=0pt of b0m] (b0t) {Data Tx: [Issuance BI]};

%--- Blok 1 ---
\node[blkhdr] (b1h)  at (4.8, 0)   {Blok 1};
\node[blkbody, below=0pt of b1h] (b1p) {Hash sebelumnya: \texttt{H(Blok\,0)}};
\node[blkbody, below=0pt of b1p] (b1m) {Merkle Root: \texttt{C3D4}};
\node[blkbody, below=0pt of b1m] (b1t) {Data Tx: [Transfer]};

%--- Blok 2 ---
\node[blkhdr] (b2h)  at (9.6, 0)   {Blok 2};
\node[blkbody, below=0pt of b2h] (b2p) {Hash sebelumnya: \texttt{H(Blok\,1)}};
\node[blkbody, below=0pt of b2p] (b2m) {Merkle Root: \texttt{E5F6}};
\node[blkbody, below=0pt of b2m] (b2t) {Data Tx: [Issuance, Transfer]};

%--- Panah antar blok ---
\draw[ar] (b0h.east |- b1p.west) -- node[above,font=\scriptsize]{hash} (b1p.west);
\draw[ar] (b1h.east |- b2p.west) -- node[above,font=\scriptsize]{hash} (b2p.west);

\node[font=\small] at (12.0, -0.6) {$\cdots$};

\end{tikzpicture}%
}
\caption{Struktur rantai blok: setiap blok memuat \textit{hash} kriptografis dari blok sebelumnya}
\label{fig:block-structure}
\end{figure}

\begin{figure}[H]
\centering
\resizebox{0.72\textwidth}{!}{%
\begin{tikzpicture}[
  leaf/.style={draw=black, fill=white, rectangle, minimum width=2.0cm,
    minimum height=0.6cm, align=center, font=\scriptsize},
  inode/.style={draw=black, fill=gray!15, rectangle, minimum width=2.4cm,
    minimum height=0.6cm, align=center, font=\scriptsize},
  root/.style={draw=black, fill=gray!35, rectangle, minimum width=2.6cm,
    minimum height=0.7cm, align=center, font=\scriptsize\bfseries},
  ar/.style={-{Stealth}, thick},
]

%--- Daun ---
\node[leaf] (t1) at (0,   0) {H(Tx\textsubscript{1})};
\node[leaf] (t2) at (2.6, 0) {H(Tx\textsubscript{2})};
\node[leaf] (t3) at (5.2, 0) {H(Tx\textsubscript{3})};
\node[leaf] (t4) at (7.8, 0) {H(Tx\textsubscript{4})};

%--- Tingkat 1 ---
\node[inode] (n12) at (1.3, 1.4) {H(T\textsubscript{1}+T\textsubscript{2})};
\node[inode] (n34) at (6.5, 1.4) {H(T\textsubscript{3}+T\textsubscript{4})};

%--- Akar ---
\node[root] (root) at (3.9, 2.8) {Merkle Root};

%--- Panah ---
\draw[ar] (t1.north)  -- (n12.south west);
\draw[ar] (t2.north)  -- (n12.south east);
\draw[ar] (t3.north)  -- (n34.south west);
\draw[ar] (t4.north)  -- (n34.south east);
\draw[ar] (n12.north) -- (root.south west);
\draw[ar] (n34.north) -- (root.south east);

\end{tikzpicture}%
}
\caption{Pohon Merkle: \textit{hash} transaksi digabungkan secara berjenjang membentuk Merkle Root yang disimpan dalam kepala blok}
\label{fig:merkle-tree}
\end{figure}


\subsubsection{Algoritme Konsensus}

Algoritme konsensus menjamin seluruh simpul pada jaringan menyepakati satu
urutan transaksi yang sama. Pada \textit{blockchain} publik digunakan mekanisme
seperti \textit{proof-of-work} (PoW) dan \textit{proof-of-stake} (PoS) yang
dirancang tahan terhadap perilaku jahat namun berbiaya komputasi tinggi.
Sebaliknya, \textit{blockchain permissioned} dapat memakai varian
\textit{Byzantine Fault Tolerance} (BFT) maupun \textit{Crash Fault Tolerance}
(CFT) seperti Raft. Identitas peserta mendukung autentikasi dan atribusi, tetapi
pemilihan BFT atau CFT tetap harus mengikuti jenis kegagalan yang diasumsikan
\cite{consensusanalysis, androulaki2018}. Algoritme Raft membagi peran simpul
menjadi tiga \textit{state}: \textit{Leader}, \textit{Candidate}, dan
\textit{Follower}. Diagram transisi \textit{state} Raft ditunjukkan pada
Gambar~\ref{fig:raft-states}.

% Diagram transisi state algoritme konsensus Raft (TikZ automata).
% Di-input dari BAB II §2.2.1.2.

\begin{figure}[H]
\centering
\resizebox{0.78\textwidth}{!}{%
\begin{tikzpicture}[
  shorten >=1pt,
  node distance=4.2cm,
  on grid,
  auto,
  state/.style={draw=black, circle, fill=white, minimum size=2.0cm,
    align=center, font=\footnotesize\bfseries},
  ar/.style={-{Stealth}, thick},
  lab/.style={font=\scriptsize, align=center},
]

\node[state] (follower)   at (0,   0)   {Follower};
\node[state] (candidate)  at (5.2, 0)   {Candidate};
\node[state] (leader)     at (10.4, 0)  {Leader};

% Follower → Candidate
\draw[ar] (follower)  to[bend left=20]
  node[above, lab]{\textit{timeout;} mulai\\\textit{election}} (candidate);

% Candidate → Leader
\draw[ar] (candidate) to[bend left=20]
  node[above, lab]{suara mayoritas\\diperoleh} (leader);

% Candidate → Follower
\draw[ar] (candidate) to[bend left=20]
  node[below, lab]{\textit{leader} ditemukan\\ atau \textit{term} lebih tinggi} (follower);

% Leader → Follower
\draw[ar] (leader) to[bend left=20]
  node[below, lab]{\textit{term} lebih tinggi\\\textit{(step down)}} (candidate);

% Self-loop Leader: heartbeat
\draw[ar] (leader.north) to[loop above, looseness=5]
  node[above, lab]{\textit{heartbeat}\\\textit{AppendEntries}} (leader.north);

% Self-loop Follower: receive heartbeat
\draw[ar] (follower.south) to[loop below, looseness=5]
  node[below, lab]{terima \textit{heartbeat};\\\textit{reset timer}} (follower.south);

\end{tikzpicture}%
}
\caption{Diagram transisi \textit{state} algoritme konsensus Raft: setiap simpul berstatus \textit{Follower}, \textit{Candidate}, atau \textit{Leader}}
\label{fig:raft-states}
\end{figure}


\subsubsection{Tipe Blockchain}

Berdasarkan model aksesnya, \textit{blockchain} dapat dikelompokkan menjadi
\textit{public} (permissionless), \textit{private}, dan \textit{consortium}
atau \textit{permissioned} \cite{bhawana2021}. Pada \textit{public blockchain},
siapa pun dapat bergabung dan berpartisipasi dalam konsensus, sehingga
memerlukan mekanisme seperti \textit{proof-of-work} yang mahal secara
komputasi. Sebaliknya, pada \textit{permissioned blockchain}, partisipasi
dibatasi pada entitas yang teridentifikasi dan memperoleh izin, sehingga
konsensus dapat menggunakan algoritme yang lebih efisien dan menghasilkan
\textit{throughput} lebih tinggi. Untuk kebutuhan CBDC, \textit{permissioned
blockchain} menjadi pilihan yang sesuai karena bank sentral memerlukan kendali
tata kelola, identitas peserta yang jelas, dan kinerja yang memadai
\cite{sethaput2025, unknown2025i}.

\subsection{Hyperledger Fabric}

\subsubsection{Arsitektur Hyperledger Fabric}

Hyperledger Fabric adalah kerangka kerja \textit{permissioned blockchain}
sumber terbuka di bawah naungan \textit{Hyperledger Foundation} yang dirancang
untuk aplikasi perusahaan \cite{androulaki2018}. Fabric memiliki arsitektur
modular dengan komponen utama sebagai berikut.

\begin{enumerate}
    \item \textbf{Peer}, yaitu simpul yang menyimpan salinan \textit{ledger} dan
    menjalankan \textit{chaincode}. \textit{Peer} berperan sebagai
    \textit{endorser} yang menyimulasikan transaksi dan sebagai
    \textit{committer} yang menulis blok ke \textit{ledger}.
    \item \textbf{Orderer}, yaitu layanan pengurut (\textit{ordering service})
    yang menyusun transaksi ke dalam blok dan mendistribusikannya ke
    \textit{peer}. Fabric mendukung konsensus berbasis Raft (etcdraft) yang
    bersifat \textit{crash fault tolerant}.
    \item \textbf{Membership Service Provider} (MSP), yaitu komponen yang
    mengelola identitas dan hak akses peserta berdasarkan sertifikat digital
    \textit{X.509} yang diterbitkan oleh \textit{Certificate Authority} (CA).
    \item \textbf{Channel}, yaitu mekanisme pemisahan jaringan yang memungkinkan
    subkelompok peserta bertransaksi secara privat dengan \textit{ledger}
    tersendiri.
    \item \textbf{Chaincode}, yaitu program \textit{smart contract} yang memuat
    logika bisnis dan aturan transaksi. Pada Fabric, \textit{chaincode} dapat
    ditulis dalam bahasa Go, Java, atau Node.js.
\end{enumerate}

Gambar~\ref{fig:fabric-arch} menunjukkan hubungan antara aplikasi klien,
\textit{endorsing peer}, \textit{ordering service}, dan \textit{channel}
dalam arsitektur Hyperledger Fabric.

% Arsitektur komponen Hyperledger Fabric (TikZ).
% Di-input dari BAB II §2.2.3.1.

\begin{figure}[H]
\centering
\resizebox{\textwidth}{!}{%
\begin{tikzpicture}[
  comp/.style={draw=black, fill=white, rectangle, rounded corners=2pt,
    minimum width=2.8cm, minimum height=0.9cm, align=center, font=\footnotesize},
  peerbox/.style={draw=black, fill=none, rectangle, rounded corners=4pt,
    inner sep=8pt},
  ordbox/.style={draw=black, fill=black!8, rectangle, rounded corners=3pt,
    minimum width=3.4cm, minimum height=1.0cm, align=center, font=\footnotesize\bfseries},
  clibox/.style={draw=black, fill=white, rectangle, rounded corners=3pt,
    minimum width=2.6cm, minimum height=0.9cm, align=center, font=\footnotesize},
  cabox/.style={draw=black, fill=gray!15, rectangle, rounded corners=3pt,
    minimum width=2.6cm, minimum height=0.9cm, align=center, font=\footnotesize},
  ledger/.style={draw=black, fill=gray!20, rectangle,
    minimum width=2.4cm, minimum height=0.6cm, align=center, font=\scriptsize},
  chan/.style={draw=black, dashed, rounded corners=5pt, fill=none, inner sep=10pt},
  ar/.style={-{Stealth}, thick},
  dar/.style={-{Stealth}, thick, dashed},
  lab/.style={font=\scriptsize, midway, align=center},
]

%--- Aplikasi Klien ---
\node[clibox] (app) at (0, 6.0) {Aplikasi Klien\\(SDK Go/Node.js)};

%--- Certificate Authority ---
\node[cabox] (ca) at (10, 6.0) {Certificate\\Authority (CA)\\MSP};

%--- Ordering Service ---
\node[ordbox] (ord) at (5, 6.0) {Ordering Service\\(Raft Leader)};

%--- Peer Org1 ---
\node[font=\footnotesize\bfseries] (p1t) at (1.5, 3.95) {Peer Org\,1};
\node[comp] (p1cc) at (1.5, 3.2) {Chaincode};
\node[comp] (p1e)  at (1.5, 2.2) {Endorser};
\node[ledger] (p1l) at (1.5, 1.2) {Ledger (WS + BC)};
\node[peerbox, fit=(p1t)(p1cc)(p1e)(p1l)] (peer1) {};

%--- Peer Org2 ---
\node[font=\footnotesize\bfseries] (p2t) at (5, 3.95) {Peer Org\,2};
\node[comp] (p2cc) at (5, 3.2) {Chaincode};
\node[comp] (p2e)  at (5, 2.2) {Committer};
\node[ledger] (p2l) at (5, 1.2) {Ledger (WS + BC)};
\node[peerbox, fit=(p2t)(p2cc)(p2e)(p2l)] (peer2) {};

%--- Peer Org3 ---
\node[font=\footnotesize\bfseries] (p3t) at (8.5, 3.95) {Peer Org\,3};
\node[comp] (p3cc) at (8.5, 3.2) {Chaincode};
\node[comp] (p3e)  at (8.5, 2.2) {Committer};
\node[ledger] (p3l) at (8.5, 1.2) {Ledger (WS + BC)};
\node[peerbox, fit=(p3t)(p3cc)(p3e)(p3l)] (peer3) {};

%--- Channel bounding box ---
\node[chan, fit=(peer1)(peer2)(peer3),
  label={[font=\footnotesize\bfseries]below:\textit{Channel} (kebijakan dukungan mayoritas)}] (chan) {};

%--- Execute-Order-Validate flow ---
\draw[ar] ([xshift=-0.45cm]app.south) to[out=-90,in=145]
  node[lab,left,pos=0.52]{Proposal transaksi}
  ([xshift=-0.45cm]peer1.north);
\draw[ar] ([xshift=0.45cm]peer1.north) to[out=35,in=-90]
  node[lab,right,pos=0.48]{RWSet + tanda tangan}
  ([xshift=0.45cm]app.south);
\draw[ar] (app.east) --
  node[lab,above]{Transaksi\\terendorsement}
  (ord.west);
\draw[ar] (ord.south) to[out=-90,in=90]
  node[lab,fill=white,inner sep=1pt,right,pos=0.15]{Blok}
  ([xshift=0.85cm]peer2.north);
\draw[ar] ([xshift=-0.75cm]ord.south) to[out=-90,in=90]
  ([xshift=0.85cm]peer1.north);
\draw[ar] ([xshift=0.75cm]ord.south) to[out=-90,in=90]
  ([xshift=0.85cm]peer3.north);
\draw[dar] (ca.south) to[out=-90,in=0]
  node[lab,right]{Sertifikat} (peer3.east);
\draw[dar] (ca.south west) to[out=-150,in=90] (ord.north east);

\end{tikzpicture}%
}
\caption{Arsitektur komponen Hyperledger Fabric: klien mengirim proposal ke \textit{peer endorser}, mengirim transaksi terendorsement ke \textit{orderer}, lalu menerima komitmen blok pada semua \textit{peer} dalam \textit{channel}.}
\label{fig:fabric-arch}
\end{figure}


\subsubsection{Hyperledger Fabric Ledger}

\textit{Ledger} pada Fabric terdiri atas dua bagian, yaitu \textit{world state}
yang menyimpan nilai terkini dari setiap kunci (umumnya menggunakan basis data
LevelDB atau CouchDB), dan \textit{blockchain} yang menyimpan riwayat transaksi
secara berurutan dalam struktur tambah-saja. Penggunaan CouchDB memungkinkan
penyimpanan dokumen JSON serta kueri kaya (\textit{rich query})
\cite{unknown2025p}. Dualitas kedua komponen \textit{ledger} tersebut
diilustrasikan pada Gambar~\ref{fig:fabric-ledger}: \textit{blockchain}
berfungsi sebagai catatan historis yang perubahan retrospektifnya dapat
dideteksi, sedangkan
\textit{world state} menyimpan nilai terkini setiap kunci yang dapat dibaca
secara langsung dan efisien oleh \textit{chaincode}.

% Dualitas komponen ledger Hyperledger Fabric (TikZ).
% Di-input dari BAB II §2.2.3.2.

\begin{figure}[H]
\centering
\resizebox{0.90\textwidth}{!}{%
\begin{tikzpicture}[
  blk/.style={draw=black, fill=gray!15, rectangle, minimum width=1.8cm,
    minimum height=1.0cm, align=center, font=\scriptsize},
  wsrow/.style={draw=black, fill=white, rectangle, minimum width=3.8cm,
    minimum height=0.5cm, align=center, font=\scriptsize},
  wshdr/.style={draw=black, fill=gray!25, rectangle, minimum width=3.8cm,
    minimum height=0.6cm, align=center, font=\scriptsize\bfseries},
  boxlabel/.style={draw=black, dashed, rounded corners=4pt, fill=none, inner sep=8pt},
  ar/.style={-{Stealth}, thick},
  lab/.style={font=\scriptsize, midway, align=center},
]

%--- Sisi kiri: Blockchain (catatan historis) ---
% Block chain
\node[blk] (bl1) at (0, 0)    {Blk 0};
\node[blk] (bl2) at (2.0, 0)  {Blk 1};
\node[blk] (bl3) at (4.0, 0)  {Blk 2};
\node[font=\small] at (5.2, 0) {$\cdots$};

\draw[ar] (bl1.east) -- (bl2.west);
\draw[ar] (bl2.east) -- (bl3.west);

\node[boxlabel, fit=(bl1)(bl2)(bl3),
  label={[font=\footnotesize\bfseries]above:\textit{Blockchain} (Catatan Historis Transaksi)},
  label={[font=\scriptsize]below:tambah-saja; perubahan dapat dideteksi}] (bcbox) {};

%--- Sisi kanan: World State (CouchDB) ---
\node[wshdr] (wsh) at (9.5, 1.0) {World State (CouchDB)};
\node[wsrow, below=0pt of wsh] (ws1) {\texttt{wallet:alice} → \{balance:500000\}};
\node[wsrow, below=0pt of ws1] (ws2) {\texttt{wallet:bob} → \{balance:200000\}};
\node[wsrow, below=0pt of ws2] (ws3) {\texttt{kyc:alice} → \{tier:STANDARD\}};
\node[wsrow, below=0pt of ws3] (ws4) {\texttt{kyc:profil001} → \{status:approved, risk:low\}};

\node[boxlabel, fit=(wsh)(ws1)(ws2)(ws3)(ws4),
  label={[font=\scriptsize]below:nilai terkini; dapat dikueri (\textit{rich query})}] (wsbox) {};

%--- Panah: Blockchain memperbarui World State ---
\draw[ar, thick] (bcbox.east) -- node[lab,above]{Validasi \&\\\textit{commit}} (wsbox.west);

\end{tikzpicture}%
}
\caption{Dualitas komponen \textit{ledger} Hyperledger Fabric: \textit{blockchain} sebagai catatan transaksi historis tambah-saja yang perubahan retrospektifnya dapat dideteksi, dan \textit{world state} sebagai representasi status terkini pada CouchDB}
\label{fig:fabric-ledger}
\end{figure}


\subsubsection{Alur Transaksi Hyperledger Fabric}

Karakteristik utama Fabric adalah model eksekusi
\textit{execute-order-validate} yang membedakannya dari model
\textit{order-execute} pada banyak \textit{blockchain} lain. Alur transaksi
Fabric berlangsung dalam lima tahap. (1) Aplikasi klien mengirim usulan
transaksi (\textit{transaction proposal}) ke sejumlah \textit{peer endorser}.
(2) \textit{Endorser} menyimulasikan transaksi terhadap \textit{chaincode} dan
mengembalikan tanda tangan dukungan (\textit{endorsement}) beserta himpunan
baca-tulis (\textit{read-write set}). (3) Klien mengumpulkan dukungan sesuai
kebijakan dukungan (\textit{endorsement policy}), lalu mengirimkannya ke
\textit{ordering service}. (4) \textit{Orderer} mengurutkan transaksi ke dalam
blok. (5) Setiap \textit{peer} memvalidasi transaksi dengan memeriksa kebijakan
dukungan dan konflik versi, kemudian menuliskannya ke \textit{ledger}
\cite{androulaki2018}. Model ini meningkatkan skalabilitas karena eksekusi
\textit{chaincode} dilakukan secara paralel sebelum pengurutan.

Pemeriksaan konflik versi pada tahap validasi menerapkan kontrol konkurensi
multiversi (\textit{Multiversion Concurrency Control}, MVCC). Setiap transaksi
membawa himpunan baca (\textit{read set}) berisi versi kunci status yang dibaca
ketika endorsemen. Sebelum \textit{commit}, \textit{peer} membandingkan versi
tersebut dengan versi terkini pada \textit{ledger}. Apabila sebuah kunci telah
berubah oleh transaksi lain dalam blok yang sama atau blok sebelumnya, versi
tidak lagi cocok sehingga transaksi ditandai tidak valid sebagai konflik MVCC
meskipun endorsemennya sah \cite{androulaki2018}. Mekanisme ini menjaga
konsistensi status pada penulisan bersamaan dan menjadi dasar analisis
perebutan dompet (\textit{hot-wallet contention}) pada BAB IV.

\textit{Endorsement policy} menentukan organisasi mana yang dukungannya
diperlukan agar suatu transaksi dianggap sah. Pada jaringan multipihak,
kebijakan seperti \textit{majority} mensyaratkan dukungan dari mayoritas
organisasi, sehingga tidak ada satu pihak pun yang dapat memutakhirkan
\textit{ledger} secara sepihak.

Gambar~\ref{fig:konsensus-raft} mengilustrasikan siklus transaksi Hyperledger
Fabric secara lengkap, mulai dari \textit{proposal} ke \textit{endorsing peer}
hingga \textit{commit} ke \textit{ledger} melalui mekanisme konsensus Raft.

% Siklus transaksi Hyperledger Fabric dengan konsensus Raft (TikZ).
% Di-input dari BAB II pada subbagian Hyperledger Fabric.

\begin{figure}[H]
\centering
\resizebox{\textwidth}{!}{%
\begin{tikzpicture}[
  proc/.style={draw=black, fill=white, rectangle, rounded corners=3pt,
    minimum width=3.4cm, minimum height=1.0cm, align=center, font=\footnotesize},
  raftbox/.style={draw=black, fill=gray!10, rectangle, rounded corners=3pt,
    minimum width=3.4cm, minimum height=1.0cm, align=center, font=\footnotesize},
  commitbox/.style={draw=black, fill=black!8, rectangle, rounded corners=3pt,
    minimum width=3.4cm, minimum height=1.0cm, align=center, font=\footnotesize},
  stepnum/.style={draw=black, circle, fill=black, text=white,
    inner sep=0pt, minimum size=0.45cm, font=\footnotesize\bfseries},
  phasebox/.style={draw=black, dashed, rounded corners=5pt,
    fill=none, inner sep=10pt},
  ar/.style={-{Stealth}, thick},
  pathlabel/.style={font=\scriptsize, midway, align=center},
]

%--- Fase 1: Endorsement ---
\node[proc]     (client1)   at (0,    0)    {Klien};
\node[proc]     (endorse)   at (0,   -2.4)  {Endorsing Peers\\(BI, DISTRIBUTOR,\\ENDUSER)};
\node[stepnum]  at (-0.15, -0.85) {1};
\node[stepnum]  at ( 0.15, -1.55) {2};

\draw[ar] ([xshift=-0.65cm]client1.south) --
  node[pathlabel,left]{\textit{proposal tx}\\(simulasi)}
  ([xshift=-0.65cm]endorse.north);
\draw[ar] ([xshift=0.65cm]endorse.north) --
  node[pathlabel,right]{RWSet +\\tanda tangan}
  ([xshift=0.65cm]client1.south);

\node[phasebox, fit=(client1)(endorse),
  label={[font=\footnotesize\bfseries]above:Fase \textit{Endorsement}}] (ph1) {};

%--- Fase 2: Ordering (Raft) ---
\node[proc]    (client2) at (5.75,  0)   {Klien\\(kumpul \textit{endorsement})};
\node[raftbox] (leader)  at (5.75, -2.9) {Orderer tunggal (Raft)\\{\scriptsize orderer.centralbank.bi-coin:7050}};
\node[stepnum] at (3.75, -1.0) {3};

\draw[ar] (client2.south) --
  node[pathlabel,right]{submit \textit{tx}\\terendorsement} (leader.north);
\node[phasebox, fit=(client2)(leader),
  label={[font=\footnotesize\bfseries]above:Fase \textit{Ordering} (Raft)}] (ph2) {};

%--- Fase 3: Validasi & Commit ---
\node[proc]      (peers)    at (11.6,  0)    {Semua Peer\\(empat organisasi)};
\node[proc]      (validate) at (11.6, -2.0)  {Validasi:\\MVCC + kebijakan\\
                                               \textit{endorsement}};
\node[commitbox] (world)    at (11.6, -3.8)  {Commit ke \textit{ledger}\\
                                               (World State / CouchDB)};
\node[stepnum] at (9.6, -1.0) {4};
\node[stepnum] at (9.6, -2.9) {5};

\draw[ar] (peers.south) --
  node[pathlabel,right]{periksa blok} (validate.north);
\draw[ar] (validate.south) --
  node[pathlabel,right]{tulis status} (world.north);

\node[phasebox, fit=(peers)(validate)(world),
  label={[font=\footnotesize\bfseries]above:Fase Validasi \& \textit{Commit}}] (ph3) {};

%--- Inter-phase arrows ---
\draw[ar] (client1.east) --
  node[pathlabel,above]{\textit{Tx} + RWSet} (client2.west);
\draw[ar] (ph2.east) -- node[pathlabel,above]{Blok terurut} (ph3.west);

\end{tikzpicture}%
}
\caption{Siklus transaksi Hyperledger Fabric pada topologi uji dengan satu \textit{consenter} Raft tanpa toleransi kegagalan \textit{orderer}.}
\label{fig:konsensus-raft}
\end{figure}


\subsection{Hyperledger Caliper}

Hyperledger Caliper adalah kerangka kerja \textit{benchmark} kinerja
\textit{blockchain} yang memungkinkan pengukuran kinerja terhadap sistem
\textit{blockchain} dengan beban kerja yang dapat dikustomisasi
\cite{hlf_caliper}. Caliper menghasilkan beban transaksi melalui sejumlah
proses pekerja (\textit{worker}) dan mengukur metrik kinerja utama berdasarkan
metrik standar \textit{Hyperledger Performance and Scale Working Group}
\cite{hlf_metrics}, antara lain:
\begin{enumerate}
    \item \textbf{Throughput}, yaitu jumlah transaksi yang berhasil diproses per
    satuan waktu, dinyatakan dalam \textit{transactions per second} (TPS).
    \item \textbf{Latency}, yaitu selang waktu antara pengiriman transaksi hingga
    transaksi tersebut tercatat dan terkonfirmasi pada \textit{ledger}.
    \item \textbf{Send rate}, yaitu laju pengiriman transaksi yang dikendalikan
    oleh \textit{rate controller}.
\end{enumerate}
Caliper juga dapat memantau penggunaan sumber daya \textit{container} seperti
penggunaan CPU dan memori selama pengujian, sehingga sesuai untuk mengevaluasi
kelayakan kinerja sistem CBDC \cite{fabric_perf2023}.

\subsection{Transfer Ritel Langsung}

Unit pembayaran pada purwarupa adalah fungsi \texttt{Transfer} yang sama untuk
transfer antarpelanggan dan pelanggan-ke-pedagang. Penyatuan jalur ini menjaga
pemeriksaan otorisasi, saldo, dan pencatatan audit tetap deterministik.
Penelitian tidak memodelkan standar instrumen pembayaran
atau proses akseptasi eksternal; fokusnya adalah kebijakan buku besar CBDC ritel.

\subsection{Application Programming Interface}

\textit{Application Programming Interface} (API) adalah antarmuka yang
memungkinkan dua perangkat lunak saling berkomunikasi melalui kontrak yang
terdefinisi. Sistem yang dibangun menggunakan \textit{REST API}
(\textit{Representational State Transfer}) \cite{fielding2000rest} sebagai antarmuka komunikasi menuju
lapisan \textit{backend} melalui protokol HTTP. API ini menjembatani permintaan
pengguna atau alat uji ke \textit{chaincode} melalui \textit{Fabric Gateway SDK}
tanpa menjadikan antarmuka pengguna sebagai bagian kontribusi penelitian.

\subsection{Bahasa Pemrograman Go}

Go adalah bahasa pemrograman terkompilasi yang dikembangkan oleh Google dengan
penekanan pada kesederhanaan dan dukungan konkurensi melalui \textit{goroutine} \cite{donovan2015go}.
Pada penelitian ini Go digunakan untuk menulis \textit{chaincode} dan
\textit{backend} karena dukungannya yang kuat terhadap Hyperledger Fabric
melalui \textit{Fabric Contract API} dan \textit{Fabric Gateway SDK}, serta
kinerja konkurensi yang baik untuk melayani banyak permintaan secara
bersamaan.

\section{Analisis Perbandingan Metode}

Dua keputusan desain utama penelitian ini dibandingkan secara sistematis, yaitu
pemilihan platform DLT dan pemilihan alat pengujian kinerja. Tabel perbandingan
memperlihatkan perbedaan antar-alternatif secara eksplisit sebagai dasar
pertanggungjawaban ilmiah.

\subsection{Perbandingan Hyperledger Fabric dengan Kerangka Kerja Blockchain Lain}
\paragraph{Perbandingan Platform DLT.}
Tabel~\ref{tab:platform-dlt} membandingkan empat platform DLT
\textit{permissioned} yang lazim digunakan dalam sistem keuangan dan CBDC:
Hyperledger Fabric, R3 Corda, Quorum/Hyperledger Besu, dan Hyperledger
Sawtooth \cite{compendium_blockchain, androulaki2018, sethaput2025}.

\begingroup
\renewcommand{\arraystretch}{1.3}
\setlength{\tabcolsep}{3pt}
\footnotesize
\begin{table}[H]
\caption{Perbandingan platform DLT \textit{permissioned} untuk sistem CBDC}
\label{tab:platform-dlt}
\centering
\begin{tabular}{>{\RaggedRight\arraybackslash}p{2.5cm}
                >{\RaggedRight\arraybackslash}p{2.8cm}
                >{\RaggedRight\arraybackslash}p{2.4cm}
                >{\RaggedRight\arraybackslash}p{2.4cm}
                >{\RaggedRight\arraybackslash}p{2.4cm}}
\toprule
\textbf{Kriteria} &
\textbf{Hyperledger Fabric} &
\textbf{R3 Corda} &
\textbf{\shortstack{Quorum/\\Besu}} &
\textbf{H. Sawtooth} \\
\midrule
Konsensus &
Raft (CFT), \textit{pluggable} &
Notary (CFT) &
IBFT/Clique &
PoET, Raft \\
Model eksekusi &
\textit{Execute-order-validate} &
\textit{Flow-based} &
\textit{Order-execute} &
\textit{Order-execute} \\
Bahasa \textit{smart contract} &
Go, Java, Node.js &
Kotlin, Java &
Solidity (EVM) &
Python, Go, Java \\
Manajemen identitas &
MSP + X.509 CA &
X.500 PKI &
Berbasis akun &
Berbasis peran \\
Privasi &
\textit{Channel}, \textit{private data} &
\textit{Need-to-know} &
TX privat &
Terbatas \\
\textit{Throughput} &
$>$3.500 TPS \cite{androulaki2018} &
${\sim}$170 TPS \cite{compendium_blockchain} &
${\sim}$300 TPS (indikatif) &
${\sim}$1.000 TPS (indikatif) \\
Penerapan CBDC yang dilaporkan &
Tidak dilaporkan &
PoC Proyek Garuda \cite{bi_garuda_poc} &
PoC Proyek Garuda \cite{bi_garuda_poc} &
Tidak \\
\bottomrule
\end{tabular}
\end{table}
\endgroup

\noindent\textbf{Simpulan:} Hyperledger Fabric dipilih karena identitas
organisasi berbasis MSP, model eksekusi modular, dukungan kontrak Go, dan
integrasi langsung dengan alat uji Caliper. MSP mengautentikasi organisasi
jaringan. \textit{Channel} dan data privat
merupakan kapabilitas platform, tetapi purwarupa penelitian memakai satu
\textit{application channel} dan tidak memisahkan buku besar \textit{wholesale}
dan ritel. PoC Proyek Garuda menggunakan Corda dan Hyperledger Besu; laporan itu
menjadi pembanding kebutuhan Indonesia, bukan bukti penerapan Fabric
\cite{bi_garuda_poc}.

\subsection{Pemilihan Alat Pengujian Blockchain}
Hyperledger Caliper dan BlockBench sama-sama dapat mengevaluasi jaringan
\textit{blockchain permissioned}, tetapi memiliki cakupan berbeda. Caliper
menjalankan beban kerja yang dikonfigurasi terhadap kontrak aktual dan
melaporkan metrik tingkat sistem, sedangkan BlockBench dirancang untuk
membandingkan kinerja keseluruhan dan komponen internal beberapa platform
\cite{hlf_caliper,dinh2017blockbench}. Perbandingan ringkas disajikan pada
Tabel~\ref{tab:caliper-blockbench}.

\begin{table}[H]
    \caption{Perbandingan Hyperledger Caliper dan BlockBench berdasarkan \cite{hlf_caliper,dinh2017blockbench}.}
    \label{tab:caliper-blockbench}
    \centering
    \small
    \begin{tabular}{@{}p{2.7cm}p{4.7cm}p{4.7cm}@{}}
        \toprule
        \textbf{Kriteria} & \textbf{Hyperledger Caliper} & \textbf{BlockBench} \\
        \midrule
        Fokus & Pengujian jaringan dengan beban kerja dan kontrak yang dapat dikonfigurasi & Perbandingan lintas platform pada tingkat makro dan mikro \\
        Metrik utama & Tingkat keberhasilan, \textit{throughput}, latensi, dan penggunaan sumber daya & \textit{Throughput}, latensi, skalabilitas, dan toleransi kegagalan \\
        Beban kerja & Modul beban kerja yang memanggil fungsi kontrak aktual & Kontrak pintar nyata dan sintetis melalui API benchmark \\
        Kesesuaian & Integrasi langsung dengan jaringan Fabric dan \textit{chaincode} penelitian & Tepat untuk analisis komponen dan perbandingan beberapa platform \\
        \bottomrule
    \end{tabular}
\end{table}

\noindent\textbf{Simpulan:} Hyperledger Caliper dipilih karena tujuan penelitian
adalah mengukur perilaku \textit{chaincode} CBDC ritel aktual pada satu jaringan
Fabric, bukan membandingkan implementasi internal beberapa platform.
